% =========================================================================
% Secao 6B - Razoes da discrepancia PNADC vs INEP
% =========================================================================

\section{Razões para a discrepância PNADC vs INEP}
\label{sec:razoes}

A análise apresentada na Seção~\ref{sec:inep} mostra que mesmo após
as calibrações metodológicas adotadas, persiste alguma divergência
entre as taxas calculadas pela PNADC longitudinal e as oficialmente
divulgadas pelo INEP. Esta seção discute, de forma sistemática, as
cinco fontes mais plausíveis para essa diferença, classificando-as
entre fontes substantivas, ou seja, refletindo medições genuinamente
distintas do fenômeno, e fontes operacionais, ou seja, decorrentes
de convenções de coleta e de definição.

A primeira fonte, de natureza operacional, é a definição de
promoção no terceiro ano do EM. Conforme detalhado na
Seção~\ref{sec:inep} e implementado na Tabela~\ref{tab:t1_brasil} via
a variável \texttt{V3014}, o INEP classifica como promoção qualquer
aluno que conclui o terceiro ano do EM, independentemente de
matrícula no ano seguinte. Em sua metodologia operacional inicial,
a PNADC exigia frequência escolar em $t+1$, o que excluía alunos
formados que ingressam no mercado de trabalho. Após a correção
introduzida nesta versão da nota, em que a conclusão é detectada via
\texttt{V3014}, o gap de promoção no EM caiu de cerca de trinta e
três pontos percentuais para aproximadamente doze pontos, conforme
documentado no Apêndice~\ref{app:variaveis} e no CHANGELOG do
projeto. Essa correção elimina a parcela definicional do gap e deixa
visíveis as fontes restantes.

A segunda fonte, de natureza substantiva, é a captura de retorno,
denotada por $R$ na decomposição da Seção~\ref{sec:inep}. A PNADC,
por entrevistar o domicílio, observa diretamente o aluno
independentemente da rede, da modalidade ou da unidade federativa em
que esteja matriculado. O Censo Escolar, em contraste, depende do
reporte da instituição de ensino e pode perder de vista alunos que
migram entre redes públicas e privadas, entre unidades federativas
ou entre modalidades, especialmente quando a rede de destino é mal
coberta pelo registro. A Figura~\ref{fig:captura_retorno} mostra
que, no ensino médio, entre três e cinco por cento dos alunos com
\texttt{V3002=1} em $t+1$ mudaram de rede ou modalidade entre $t$ e
$t+1$, o que estabelece um limite superior para a parcela do gap
atribuível à captura de retorno. Esses alunos são classificados
como evadidos pelo Censo, mas continuam estudando segundo a PNADC.
A magnitude desse componente é portanto da ordem de três a cinco
pontos percentuais do gap remanescente.

A terceira fonte, de natureza substantiva, é a diferença de
universo. A PNADC mede uma amostra probabilística de domicílios e
expande para a população residente, enquanto o Censo Escolar mede o
universo de matrículas declaradas pelas instituições. A diferença
entre os denominadores das duas fontes é tipicamente de dois a três
por cento, com a PNADC reportando ligeiramente menos alunos do que
o Censo. Parte dessa diferença reflete sub-cobertura do registro
administrativo, especialmente em redes pequenas e em escolas
técnicas isoladas, mas parte também reflete super-cobertura, quando
um mesmo aluno aparece em duas matrículas no Censo, fenômeno que o
INEP tenta corrigir via a chave CPF mas que não é eliminado
completamente.

A quarta fonte é o conceito temporal de frequência. A PNADC mede a
frequência na semana de referência, isto é, na semana em que a
entrevista foi conduzida. O Censo Escolar mede a matrícula
registrada para o ano letivo. Os dois conceitos divergem em
situações de ausência temporária, isto é, aluno matriculado mas
ausente na semana de referência, e em situações de matrícula
encerrada antes do final do ano letivo. A magnitude desse
componente é difícil de quantificar diretamente, mas estimativas
indiretas sugerem que está na ordem de um a três por cento.

A quinta fonte é o conceito de auto-relato versus registro
administrativo. A PNADC pergunta ao responsável pelo domicílio
sobre a situação escolar de cada morador, e a resposta pode estar
sujeita a erro de declaração, em particular para adolescentes que
estão fora do convívio diário do informante, ou para casos em que
o aluno frequenta uma instituição não declarada por motivos
diversos. O Censo Escolar, em contraste, registra a matrícula
administrativa pela instituição, com menor erro de declaração para
o aluno em si, mas possivelmente com viés institucional, em que a
instituição declara matrículas inativas para manter financiamento.
Os dois métodos têm portanto seus próprios vieses, e a comparação
entre eles é informativa precisamente porque cada um detecta o que
o outro perde de vista.

A soma quantitativa dos cinco componentes acima, para a transição
2019/2020 no ensino médio, totaliza aproximadamente doze pontos
percentuais, o que corresponde ao gap residual observado na
Tabela~\ref{tab:t4_pnadc_inep} após a correção definicional. Isto
é, a definição de promoção sozinha explica cerca de vinte pontos
do gap original de trinta e dois, e os doze pontos restantes
distribuem-se entre captura de retorno (até cinco pontos),
diferença de universo (dois a três pontos), conceito temporal (um
a três pontos), e auto-relato versus administrativo (um a dois
pontos), em ordens de magnitude consistentes com a literatura
metodológica brasileira.

A implicação prática é a seguinte. A PNADC e o INEP, após
calibradas as definições, são fontes complementares e mutuamente
informativas. A PNADC captura aspectos que o INEP perde, em
particular a captura de retorno entre redes e a heterogeneidade
socioeconômica do fluxo. O INEP captura aspectos que a PNADC perde
ou estima com ruído, em particular o conceito administrativo de
matrícula ativa em escolas pequenas e isoladas. A convergência
entre as duas fontes, após calibração, fortalece a confiança em
ambas, enquanto a divergência residual fornece informação
substantiva sobre onde cada fonte tem viés conhecido.
